\section{Experimental Setup}
\label{sec:setup}

\textbf{Benchmark.} We rebuilt the Phase~1 benchmark from the base paper's own data. VERINA-BASIC has 108 tasks, all of which the VERINA authors translated into Lean from Dafny originals in CloverBench (59) and MBPP-DFY (49). We kept every task whose types have a direct counterpart in both Dafny and Python (integers, naturals, booleans, characters, strings, lists and tuples), which left 97 tasks after removing ten with floats, options, maps, bytes, function arguments or nested arrays, and one malformed record. For each remaining task we wrote a Dafny specification, a Python/iContract specification and a reference solution in each language, using VERINA's Lean pre- and postcondition as the ground truth for what the task means and VERINA's positive tests as our tests (2 to 12 per task, median 5). A task was admitted only if, in both languages, (i) the reference passes the integrity guard, (ii) it verifies, (iii) it passes all tests, (iv) the empty stub is rejected, and (v) the specification accepts the expected output and rejects every known-wrong output that VERINA lists for each test input. We call check (v) a \emph{spec probe}. In Dafny, a probe is a copy of the method with the test input pinned by extra \texttt{requires} clauses and a body that returns the candidate output. Because both specifications answer the same concrete questions, the probes are our main evidence that the two versions of a task are equivalent. As a second check we ran both references on random inputs that satisfy the precondition and compared their outputs (differential testing). All 97 tasks passed the five checks in both languages; for three Dafny tasks the verifier could not decide some accept-probes on literal inputs, and these exemptions are documented in the task files. Differential testing on \R{eq.inputs} random inputs found no disagreement between the Dafny and Python references.

\textbf{Splits.} Before any experiment we drew 20 tasks at random (seed 2026) as a development split and kept the remaining 77 as the test split. All design decisions (five development iterations, recorded in the repository's development log), the prompts, and the two worked examples used by E2 come from the development split only. Every number in Sections~\ref{sec:results}--\ref{sec:failures} is on the test split, which gives 154 task instances (77 per language).

\textbf{Model and decoding.} All experiments use \texttt{gpt-4o-mini-2024-07-18}, one of the models VERINA evaluates, which makes the base-paper comparison like for like. We had planned to add GPT-4.1, but the project had a hard budget of 2~USD for API calls, so the stronger model could not be run at the scale we needed (Section~\ref{sec:limits}). Agents use temperature 0.2 (0.6 for rewrites), at most 4096 output tokens, OpenAI seeds 0, 1 and 2 for the three repetitions, and a budget of $K = 5$ candidates. Every request and response is stored in a cache keyed on the full request, so all runs can be replayed exactly without new API calls.

\textbf{Experiments.} E1 runs the Phase~1 baselines (single-shot and testing-only). E2 reproduces VERINA's CodeGen protocol on our tasks, with the same instruction structure (task description plus a template that contains the reference specification), two worked examples, temperature 1.0 and plain-text answers. Because of the budget we drew 3 samples per task instead of VERINA's 5, but pass@1 is still estimated without bias as the mean over samples~\cite{chen2021codex}. E3 runs the full SpecLoop agent, and E4 adds the verifier-only arm and reuses E1 and E3 for the gap-wise tests. In E5 we remove one component at a time from SpecLoop. Finally, E6 collects cost, latency, tool usage and reliability measures from all runs, together with a stress test of the integrity guard.

\textbf{Implementation.} SpecLoop is about 4,200 lines of Python. Experiments ran on a 16-core Windows~10 machine with Dafny~4.11.0, CrossHair~0.0.113 and Python~3.12, with up to 6 runs in parallel per experiment. The code, task files, raw results, trajectories and the scripts that produce every table and figure are in the project repository.

\section{Evaluation Methodology}
\label{sec:eval}

\textbf{Metrics.} Our primary metric is \emph{correctness}, the share of runs whose final program satisfies Eq.~(1). We also report the \emph{verified rate} (the verifier accepts the final program), the \emph{test pass rate} on all tests, VERINA's CodeGen metric (pass@1 on tests, with the share of samples that do not compile), and the \emph{spec-intact rate}. For reliability we report the \emph{false-success rate}, the share of runs where the agent stopped because its own checks passed but the oracle disagrees (possible when a testing-only agent never sees the verifier, or when the in-loop guard is off). For efficiency we count model calls, tokens, cost in USD, verifier and test calls, and wall-clock time per run. Per-candidate oracle verdicts give \emph{correct@$k$}, the correctness an agent would reach with a budget of $k$ candidates.

\textbf{Statistics.} Each configuration is run three times, and we report the mean and standard deviation over the three runs. For comparisons between configurations we pair results by task: we average each task's outcome over seeds and apply a Wilcoxon signed-rank test~\cite{wilcoxon1945individual} to the per-task differences, and as a second check an exact McNemar test~\cite{mcnemar1947note} on matched (task, seed) outcomes. We also give 95\% bootstrap confidence intervals for the difference in rates. For Gap~1 the pairing is by VERINA task across the two languages. We use $\alpha = 0.05$ and do not correct for multiple comparisons, so p-values near the threshold should be read with care.

\textbf{Base-paper comparison.} We compare against VERINA in two ways and keep them apart in every table. \emph{Reported} numbers are VERINA's own GPT-4o-mini results on VERINA-BASIC in Lean (Table~9 of~\cite{ye2026verina}: 44.0\% pass@1, 54.9\% of samples do not compile). \emph{Reproduced} numbers come from E2, which runs VERINA's protocol with the same model on our Dafny and Python versions of the same tasks. The difference between the two isolates the effect of the specification language, and the difference between E2 and our agents isolates the effect of the feedback loop.
